% Источники: exam/materials/преподаватель/2020/23-03-2020-MODEL L6(2).doc; exam/materials/моделирование.pdf, разделы 21--22.
\section{Анализ временных рядов в математическом моделировании.}

\textbf{Временной ряд} — совокупность последовательных измерений значений переменной (процесса), произведенных через определенные (чаще равные) значения параметра (обычно времени).

Если измеряемые значения являются многомерными, то временной ряд также называется многомерным. Для анализа полученных таким образом данных используется область статистики, называемая \textbf{анализом временных рядов}.

\subsection*{Задача прогнозирования}

\textbf{Цель:} на основе исторических данных предсказать будущие значения временного ряда.

Для решения обычно сводится к задаче регрессии. При этом важно, что для построения модели можно использовать только признаки (факторы), полученные до текущего момента времени (не из будущего).

\textbf{Факторы, которые можно считать:}
\begin{itemize}
  \item предыдущие значения целевой переменной;
  \item статистики по целевой переменной за последние $n$ отрезков времени (среднее, мин, макс за 1 месяц, 3 месяца, 180 дней и т.д.);
  \item временные признаки — время суток, сезон, выходной/будни и т.д. (если это можно посчитать);
  \item дополнительные факторы, известные из условий.
\end{itemize}

\subsection*{Смешанное прогнозирование}

Под смешанным прогнозированием обычно понимают использование комбинированного подхода, который учитывает как детерминированные (например, тренд и сезонность), так и стохастические (например, случайные колебания) элементы временного ряда.

Целью является построение модели, которая будет:
\begin{itemize}
  \item учитывать долгосрочные тенденции (тренд);
  \item улавливать периодические колебания (сезонность);
  \item корректно обрабатывать непредсказуемые шумовые компоненты.
\end{itemize}

Это особенно важно, если данные являются нестационарными, то есть их статистические характеристики со временем меняются. Примером метода, способного работать с такими данными, является SARIMA, который учитывает тренд, сезонность и стохастические колебания одновременно.

Каждую из компонент моделируют отдельно, затем возвращаются по обратной формуле к исходному ряду. Пример: по тренду можно построить линейную регрессию; под сезонность подогнать периодическую функцию типа синуса; остаточная компонента содержит случайные колебания, которые не объясняются ни трендом, ни сезонностью. В идеале эти остатки должны быть стационарными и независимыми. На них можно применить процессы ARMA$(p,q)$.

\clearpage
